%%%%%%%%%%%%%%%%%%%%%%%%%%%%%%%%%%%%%%%%%%%%%%%%%%%%%%%%%%%%%%%%%
% PHIẾU HỌC TẬP: CRYPTOCODE
%%%%%%%%%%%%%%%%%%%%%%%%%%%%%%%%%%%%%%%%%%%%%%%%%%%%%%%%%%%%%%%%%
\documentclass[12pt,a4paper]{article}

\usepackage[utf8]{inputenc}
\usepackage[T1]{fontenc}
\usepackage[vietnamese]{babel}
\usepackage{amsmath,amssymb}
\usepackage{geometry}
\usepackage{tabularx, array}
\usepackage[svgnames]{xcolor}
\usepackage{tikz}

\geometry{left=2cm,right=2cm,top=2cm,bottom=2cm}
\definecolor{cryptogreen}{RGB}{16, 185, 129}

\begin{document}

\begin{center}
    \begin{tikzpicture}
        \node[draw=cryptogreen, line width=2pt, rounded corners=10pt, inner sep=10pt] (box) {
            \begin{minipage}{0.9\textwidth}
                \centering
                \bfseries\Large\color{cryptogreen} NHIỆM VỤ ĐIỆP VIÊN: MÃ HÓA \& GIẢI MÃ
            \end{minipage}
        };
    \end{tikzpicture}
    \vspace{0.5cm}
\end{center}

\noindent \textbf{Điệp viên (Học sinh):} \dotfill \textbf{Mã số (Lớp):} \dotfill

\section*{\color{cryptogreen} BẢNG THAM CHIẾU}
\begin{center}
\begin{tabular}{|c|c|c|c|c|c|c|c|c|c|c|c|c|}
\hline
A & B & C & D & E & F & G & H & I & J & K & L & M \\
\hline
0 & 1 & 2 & 3 & 4 & 5 & 6 & 7 & 8 & 9 & 10 & 11 & 12 \\
\hline
\end{tabular}
\vspace{0.3cm}

\begin{tabular}{|c|c|c|c|c|c|c|c|c|c|c|c|c|}
\hline
N & O & P & Q & R & S & T & U & V & W & X & Y & Z \\
\hline
13 & 14 & 15 & 16 & 17 & 18 & 19 & 20 & 21 & 22 & 23 & 24 & 25 \\
\hline
\end{tabular}
\end{center}

\section*{\color{cryptogreen} NHIỆM VỤ 1: MÃ HÓA TIN NHẮN}
\textbf{Tin gốc:} HELLO WORLD \quad \textbf{Khóa:} $k = 5$

Công thức: $E(x) = (x + k) \mod 26$

\begin{center}
\begin{tabular}{|c|c|c|c|c|c|c|c|c|c|c|}
\hline
Chữ gốc & H & E & L & L & O & W & O & R & L & D \\
\hline
Giá trị $x$ & 7 & 4 & 11 & 11 & 14 & 22 & 14 & 17 & 11 & 3 \\
\hline
$x + k$ & 12 & & & & & & & & & \\
\hline
$(x+k) \mod 26$ & 12 & & & & & & & & & \\
\hline
Chữ mã hóa & M & & & & & & & & & \\
\hline
\end{tabular}
\end{center}
\textbf{Tin đã mã hóa:} M \_ \_ \_ \_ \quad \_ \_ \_ \_ \_

\section*{\color{cryptogreen} NHIỆM VỤ 2: GIẢI MÃ TIN MẬT}
\textbf{Tin mã hóa:} KHOOR ZRUOG \quad \textbf{Khóa:} $k = 3$

Công thức: $D(y) = (y - k) \mod 26$ (nếu âm thì cộng thêm 26)

\begin{center}
\begin{tabular}{|c|c|c|c|c|c|c|c|c|c|c|}
\hline
Chữ mã & K & H & O & O & R & Z & R & U & O & G \\
\hline
Giá trị $y$ & 10 & 7 & & & & & & & & \\
\hline
$y - k$ & 7 & 4 & & & & & & & & \\
\hline
$(y-k) \mod 26$ & 7 & 4 & & & & & & & & \\
\hline
Chữ gốc & H & E & & & & & & & & \\
\hline
\end{tabular}
\end{center}
\textbf{Tin giải mã:} HE \_ \_ \_ \quad \_ \_ \_ \_ \_

\section*{\color{cryptogreen} NHIỆM VỤ 3: PHÁ MÃ (Brute-force)}
\textbf{Tin mã hóa:} LIPPS \quad \textbf{Khóa:} KHÔNG BIẾT!

Thử các giá trị $k$ từ 1 đến 5 và tìm từ có nghĩa:
\begin{itemize}
    \item $k=1$: KHOOR (Vô nghĩa)
    \item $k=2$: ......... 
    \item $k=3$: .........
    \item $k=4$: .........
    \item $k=5$: .........
\end{itemize}
\textbf{Khóa đúng là:} $k = $ ......... \quad \textbf{Tin gốc:} .........

\end{document}
